% El título anterior era «Fundamentos de computación cuántica». Ya no vale: el capítulo
% cubre ahora también los modelos de aprendizaje automático. De paso desaparece el
% Overfull \hbox de 3,35 pt que obligaba a partir el título a mano.
\chapter{Fundamentos}
\label{cap:fundamentos}

% =====================================================================================
% GUION. Capitulo de CONTEXTO: recoge conocimiento que YA EXISTE y se cita.
% Reestructurado en ago-2026 en dos bloques, 2.1 cuantico y 2.2 aprendizaje automatico.
%
% 🔴 LA FRONTERA CON EL CAPITULO 9, que es la regla que evita duplicar:
%
%      2.2  → QUE ES el metodo. Conocimiento de la literatura. Se CITA.
%      cap. 9 → POR QUE ese metodo aqui, como se configura, que se midio.
%               Decisiones propias. Se JUSTIFICA con nuestros datos.
%
%    Ejemplo: «que es la regularizacion L2» va en 2.2; «por que Ridge y no minimos
%    cuadrados, dada la colinealidad extrema que medimos» va en el 9. Si una frase
%    contiene una cifra nuestra, esta en el capitulo equivocado.
%
% 🔴 TONO, y es asimetrico a proposito. El tribunal es de INTELIGENCIA ARTIFICIAL:
%      2.1 → no saben cuantica. Claro, ameno, con intuicion antes que formalismo.
%            El tecnicismo entra, pero siempre despues de haberlo hecho comprensible.
%      2.2 → hablas su idioma. Denso y rapido en Ridge y Random Forest, que conocen
%            mejor que nadie. Detenerse SOLO en la red de grafos, que es lo que
%            probablemente no han visto aplicado a esto.
%
% ⚠️ Reglamento art. 11.4.a: «no es original el trabajo que contenga un exceso de
%    informacion de otras fuentes, AUNQUE ESTEN DEBIDAMENTE CITADAS». En un capitulo
%    de fundamentos es donde mas facil es caer. Defensa: incluir solo lo que los
%    capitulos 3, 6 y 9 van a usar de verdad, y que cada apartado apunte hacia adelante.
%
% ⚠️ Solo se numera hasta \subsection (secnumdepth = 2 en la clase book). Un
%    \subsubsection saldria SIN numero. Por eso el cuarto nivel va como parrafo con
%    entradilla en negrita, igual que en el resto de la memoria.
% =====================================================================================

% Parrafo de entrada del capitulo: que hace falta saber para leer el resto y por que
% este capitulo tiene dos mitades que parecen inconexas y no lo son.
% ⚠️ La norma prohibe dos encabezados consecutivos sin texto entre medias.


% =====================================================================================
\section{Fundamentos de computación cuántica}
\label{sec:fund_cuantica}
% =====================================================================================
% Fuente principal: Nielsen y Chuang (2010). Citar con paginas concretas:
%   \cite[pp.~13--16]{nielsen2010quantum} · \citeA[cap.~4]{nielsen2010quantum}
% UNA sola entrada en bibliografia.bib para todo el libro; la pagina va en la cita.

% Parrafo de entrada de la seccion.

\subsection{Del bit al qubit}
\label{sec:bit_qubit}
% LA PUERTA DE ENTRADA. Si el tribunal se pierde aqui, se pierde en todo el capitulo.
% Arrancar por lo que SI conocen: el bit, y de ahi la diferencia.
%
% Fuente: Nielsen y Chuang (2010), seccion 1.2, p. 13.
% ⚠️ VERIFICAR LA PAGINA del ultimo parrafo (el estado |+> y el ejemplo del atomo
%    iluminado): puede estar en la p. 14 y no en la 13. No citar p. 13 para eso sin
%    comprobarlo en el libro.
El bit es el concepto fundamental de la computación y la información clásicas. La
computación y la información cuánticas se construyen sobre un concepto análogo, el bit
cuántico o, abreviadamente, el \textbf{qubit} \cite[secc.~1.2, p.~13]{nielsen2010quantum}.

¿Qué es entonces un qubit? Igual que un bit clásico tiene un estado —o 0, o 1—, un qubit
también tiene un estado. Dos estados posibles para un qubit son $\ket{0}$ y $\ket{1}$, que,
como cabe suponer, se corresponden con los estados 0 y 1 de un bit clásico. La notación
$\ket{\ }$ se denomina \textit{notación de Dirac} y es la
notación estándar para los estados en mecánica cuántica. La diferencia entre bits y qubits
es que un qubit puede encontrarse en un estado distinto de $\ket{0}$ o $\ket{1}$: también es
posible formar combinaciones lineales de estados, llamadas habitualmente
\textbf{superposiciones}:
\begin{equation}
\label{eq:superposicion}
\ket{\psi} = \alpha \ket{0} + \beta \ket{1}.
\end{equation}
Los números $\alpha$ y $\beta$ son números complejos.

Esa misma información se puede escribir de forma más simple. Como el estado queda determinado
por sus dos amplitudes, basta con apilarlas para obtener una \textbf{representación
vectorial} del qubit, tomando $\ket{0}$ y $\ket{1}$ como los dos vectores de la base:
\begin{equation}
\label{eq:vector_estado}
\ket{0} \equiv \begin{pmatrix} 1 \\ 0 \end{pmatrix},
\qquad
\ket{1} \equiv \begin{pmatrix} 0 \\ 1 \end{pmatrix},
\qquad
\ket{\psi} = \alpha\ket{0} + \beta\ket{1} \equiv
\begin{pmatrix} \alpha \\ \beta \end{pmatrix}.
\end{equation}
Las dos notaciones son equivalentes y a lo largo del documento se emplean indistintamente:
la de Dirac resulta más cómoda para razonar, y la vectorial para calcular, porque convierte
las operaciones sobre qubits en productos de matrices —como se verá en el
apartado~\ref{sec:puertas}—.

Un bit se puede examinar para determinar si está en el estado 0 o en el 1. Los ordenadores
lo hacen constantemente, por ejemplo cuando recuperan el contenido de su memoria. Cabe recalcar
que un qubit no se pueda examinar para determinar su estado cuántico, es decir, los
valores de $\alpha$ y $\beta$. En su lugar, la mecánica cuántica nos dice que solo podemos
adquirir una información mucho más restringida sobre el estado cuántico. Al medir un qubit
obtenemos o bien el resultado $0$, con probabilidad $|\alpha|^2$, o bien el resultado 1, con
probabilidad $|\beta|^2$. Naturalmente, $|\alpha|^2 + |\beta|^2 = 1$, puesto que las
probabilidades deben sumar uno \cite[p.~13]{nielsen2010quantum}.

Que un qubit pueda estar en un estado de superposición choca con nuestra comprensión de
«sentido común» del mundo físico que nos rodea. Un bit clásico es como una moneda: o cara,
o cruz. En monedas imperfectas puede haber estados intermedios, como quedarse en equilibrio
sobre el canto, pero en el caso ideal se pueden descartar. Un qubit, por el contrario, puede
existir en un continuo de estados entre $\ket{0}$ y $\ket{1}$, hasta que se lo observa
\cite[p.~13]{nielsen2010quantum}.

Aunque los qubits parezcan «raros» o contraintuitivos, son algo físico y real, no una idea
abstracta. Se han comprobado en muchos experimentos y existen distintas formas de
construirlos en el mundo real:
\begin{itemize}
    \item un qubit puede ser la polarización de un fotón, esto es, dos formas distintas de
          vibrar la luz;
    \item puede ser la orientación del espín de un núcleo dentro de un campo magnético;
    \item puede ser dos estados de un electrón en un átomo: el estado «fundamental»
          ($\ket{0}$) y el estado «excitado» ($\ket{1}$).
\end{itemize}

En este último ejemplo, el átomo funciona como un pequeño sistema de dos niveles. Si se le
ilumina con luz de la energía adecuada, se puede hacer que el electrón pase de $\ket{0}$ a
$\ket{1}$, o de $\ket{1}$ a $\ket{0}$. Hasta aquí, todo es clásico: dos estados, uno u otro.

Pero lo interesante llega después. Si no se le da suficiente tiempo de luz, el electrón no
llega completamente a $\ket{1}$, sino que se queda en un estado intermedio, una superposición
entre $\ket{0}$ y $\ket{1}$. Ese estado intermedio es lo que se denomina $\ket{+}$.

% PENDIENTE en esta misma subseccion, aun sin redactar:
%   - Definir |+> explicitamente, si se decide dar la formula: (|0> + |1>)/raiz(2).
%   - El entrelazamiento como segunda propiedad que da la ventaja computacional.
%   - Por que n qubits necesitan 2^n amplitudes, y por que eso es a la vez el argumento
%     que hace interesante el campo y el que limita el tamaño de nuestro conjunto de
%     datos (capitulo~\ref{cap:rigor}).
%   - La contrapartida que enlaza con el capitulo~\ref{cap:ruido}: el estado cuantico es
%     fragil, y esa fragilidad es el problema que este trabajo intenta mitigar.

\subsection{La esfera de Bloch y la información que guarda un qubit}
\label{sec:bloch}
% Fuente: Nielsen y Chuang (2010), seccion 1.2.
% ⚠️ VERIFICAR LAS PAGINAS: la esfera de Bloch y la figura 1.3 estan en la seccion 1.2,
%    pero NO en la p. 13 —caen algo despues, hacia las pp. 14-16—. Se cita por seccion
%    hasta comprobar el numero exacto en el libro. No inventar la pagina.
Una imagen útil para pensar en los qubits es la siguiente representación geométrica. Como
$|\alpha|^2 + |\beta|^2 = 1$, la ecuación~\eqref{eq:superposicion} se puede reescribir como
\begin{equation}
\label{eq:bloch_gamma}
\ket{\psi} = e^{i\gamma}\left( \cos\frac{\theta}{2}\ket{0}
             + e^{i\varphi} \sin\frac{\theta}{2}\ket{1} \right),
\end{equation}
donde $\theta$, $\varphi$ y $\gamma$ son números reales. El factor $e^{i\gamma}$ que aparece
delante se puede ignorar, porque no tiene efectos observables, y por esa razón se puede
escribir
\begin{equation}
\label{eq:bloch}
\ket{\psi} = \cos\frac{\theta}{2}\ket{0} + e^{i\varphi} \sin\frac{\theta}{2}\ket{1}.
\end{equation}

Los números $\theta$ y $\varphi$ definen un punto sobre la esfera unidad tridimensional,
tal como se muestra en la Figura~\ref{fig:bloch}. Esta esfera se denomina habitualmente
\textbf{esfera de Bloch}, y nos permite visualizar el estado de un solo
qubit. Muchas de las operaciones sobre qubits individuales se describen 
dentro de esta representación. Sin embargo, para varios qubits presenta una limitación
porque no se conoce ninguna generalización sencilla en la
esfera de Bloch \cite[secc.~1.2]{nielsen2010quantum}.

% ⚠️ FORMATO EXIGIDO por las Instrucciones, apartado 1.3, PAGINA 6:
%      «Titulo en la parte superior de la tabla o figura.»
%      «Fuente de la tabla o figura en la parte inferior.»
%    Por eso el \caption va ANTES del dibujo. La plantilla oficial lo hace al reves en
%    su ejemplo; aqui manda la norma escrita.
%
% FIGURA PROPIA, dibujada en TikZ. El codigo esta en figuras/cap02/bloch_tikz.tex,
% con sus parametros comentados (radio, perspectiva, y los angulos theta y phi del
% estado que se dibuja).
%
% Al ser de elaboracion propia NO arrastra ninguna licencia, y ademas usa la
% tipografia del documento, asi que las etiquetas casan con el resto del texto.
%
% ALTERNATIVA: la version descargada de Wikimedia sigue en figuras/cap02/bloch.pdf.
% Para volver a ella, sustituir el \input por
%     \includegraphics[width=0.42\textwidth]{cap02/bloch}
% y devolver el pie a: «Fuente: Glosser.ca (2012), Bloch Sphere.svg, Wikimedia
% Commons. Licencia CC BY-SA 3.0.»
\begin{figure}[htbp]
    \centering
    \caption{Un bit clásico frente a un qubit: dos estados posibles frente
             a estados infinitos sobre la esfera de Bloch.}
    \label{fig:bloch}
    % =====================================================================================
% Esfera de Bloch — figura propia en TikZ.
%
% Se incluye desde capitulos/02_fundamentos.tex con \input. Es vectorial, usa la
% tipografia del documento y es de ELABORACION PROPIA: no arrastra ninguna licencia.
%
% GEOMETRIA. El plano ecuatorial se dibuja como una elipse alineada con los ejes de
% la pagina, de semiejes \R y \K*\R. Un punto del ecuador se parametriza con el angulo
% de pantalla t:            (\R cos t , \K \R sin t)
% Los ejes se colocan sobre esa misma elipse, asi que todo queda consistente:
%     \Ty  = 0°    -> el eje y, hacia la derecha
%     \Tx  = 235°  -> el eje x, hacia abajo a la izquierda
% El angulo azimutal phi se mide DESDE el eje x hacia el eje y, luego t = \Tx + phi.
%
% PARAMETROS AJUSTABLES: \R (radio), \K (aplastamiento de la perspectiva),
% \Theta y \Phi (posicion del estado dibujado).
% =====================================================================================
% ⚠️ \shorthandoff{<>} ES NECESARIO. `babel` con la opcion `spanish` convierte < y >
%    en caracteres ACTIVOS (los usa para las comillas << >>), y eso rompe la sintaxis
%    de puntas de flecha de TikZ: «>={Stealth}» y «[->]» fallan con
%    «Argument of \language@active@arg> has an extra }».
%    Va dentro de un grupo, asi que los catcodes se restauran solos al cerrar.
\begingroup
\shorthandoff{<>}%
\begin{tikzpicture}[>={Stealth[length=2.2mm]}, line width=0.5pt]

  \pgfmathsetmacro{\R}{2.2}      % radio de la esfera
  \pgfmathsetmacro{\K}{0.30}     % achatamiento del ecuador en perspectiva
  \pgfmathsetmacro{\Tx}{235}     % angulo de pantalla del eje x
  \pgfmathsetmacro{\Theta}{40}   % angulo polar del estado dibujado
  \pgfmathsetmacro{\Phi}{55}     % angulo azimutal del estado dibujado

  % --- coordenadas del estado |psi> ------------------------------------------------
  \pgfmathsetmacro{\Tp}{\Tx+\Phi}                       % angulo de pantalla del azimut
  \pgfmathsetmacro{\Rxy}{\R*sin(\Theta)}                % radio de la proyeccion ecuatorial
  \pgfmathsetmacro{\Px}{\Rxy*cos(\Tp)}
  \pgfmathsetmacro{\Pyeq}{\K*\Rxy*sin(\Tp)}             % altura de la proyeccion
  \pgfmathsetmacro{\Py}{\Pyeq+\R*cos(\Theta)}           % altura del propio estado

  % =====================================================================================
  % EL BIT CLASICO, a la izquierda, para la comparativa.
  % ⚠️ NO TOCA NADA DE LA ESFERA: se dibuja en su propia columna, a la izquierda del
  %    origen. La esfera sigue centrada en (0,0) exactamente como estaba.
  %
  % Los dos puntos van al MISMO azul que los polos de la esfera, y alineados con ellos.
  % Es deliberado: asi se lee de un vistazo que el bit clasico tiene DOS estados y que
  % el qubit tiene esos mismos dos... y toda la superficie que hay entre medias.
  % =====================================================================================
  \pgfmathsetmacro{\Xbit}{-2.60*\R}     % columna del bit clasico
  % ⚠️ Los cuatro puntos van al MISMO radio (\Rpunto) y las dos parejas a la MISMA
  %    altura, y=+\R y y=-\R: asi el 0 queda alineado con |0> y el 1 con |1>.
  \pgfmathsetmacro{\Rpunto}{3.2}        % radio de los puntos, en pt

  \fill[rojoclasico] (\Xbit,\R) circle (\Rpunto pt);
  \node[right=5pt] at (\Xbit,\R) {$0$};
  \fill[rojoclasico] (\Xbit,-\R) circle (\Rpunto pt);
  \node[right=5pt] at (\Xbit,-\R) {$1$};

  % rotulos de cada mitad
  \node[font=\bfseries] at (\Xbit,{-1.62*\R}) {Bit clásico};
  \node[font=\bfseries] at (0,{-1.62*\R}) {Qubit};

  % --- esfera y ecuador -------------------------------------------------------------
  \draw (0,0) circle (\R);
  \draw[dashed] (0,0) ellipse ({\R} and {\K*\R});

  % --- ejes -------------------------------------------------------------------------
  % Las puntas de flecha MUEREN sobre la propia esfera, sin sobresalir; las etiquetas
  % van fuera. El eje x es mas corto a proposito: no apunta hacia el borde del circulo
  % sino hacia el punto del ECUADOR que le corresponde, que en perspectiva cae dentro.
  % eje z: atraviesa la esfera de polo a polo, con punta arriba
  \draw (0,-\R) -- (0,0);
  \draw[->] (0,0) -- (0,\R);
  % El eje z lleva SU PROPIA etiqueta en negro, como x e y. Las dos etiquetas azules
  % de arriba y abajo no nombran el eje: nombran los dos ESTADOS de la base.
  \node[right=2pt] at ({0.02*\R},{0.86*\R}) {$\hat{\mathbf{z}}$};
  % Solo el estado: el eje ya lo nombra la etiqueta negra de arriba. Centradas sobre
  % el polo, que es lo que marcan.
  \node[above=2pt, azulunir] at (0,\R) {$\ket{0}$};
  \node[below=2pt, azulunir] at (0,-\R) {$\ket{1}$};
  % eje y: hacia la derecha, hasta el borde
  \draw[->] (0,0) -- (\R,0);
  \node[right=2pt] at (\R,0) {$\hat{\mathbf{y}}$};
  % eje x: hasta su punto del ecuador, abajo a la izquierda
  \draw[->] (0,0) -- ({\R*cos(\Tx)},{\K*\R*sin(\Tx)});
  \node[below left, inner sep=1pt] at ({\R*cos(\Tx)},{\K*\R*sin(\Tx)})
        {$\hat{\mathbf{x}}$};

  % --- los dos estados de la base computacional --------------------------------------
  % |0> y |1> son EXACTAMENTE los polos de la esfera. Marcarlos es correcto y ayuda:
  % deja ver de un vistazo que los estados clasicos son solo dos puntos de todos los
  % posibles. En azul institucional, el mismo de los titulos de capitulo.
  \fill[azulunir] (0,\R) circle (\Rpunto pt);
  \fill[azulunir] (0,-\R) circle (\Rpunto pt);

  % --- proyeccion del estado sobre el plano ecuatorial -------------------------------
  \draw[dotted] (0,0) -- (\Px,\Pyeq);
  \draw[dotted] (\Px,\Pyeq) -- (\Px,\Py);

  % --- el estado --------------------------------------------------------------------
  \draw[->] (0,0) -- (\Px,\Py);
  \fill (\Px,\Py) circle (1.5pt);
  \node[above right, inner sep=2pt] at (\Px,\Py) {$\ket{\psi}$};

  % --- angulo theta: entre el eje z y el estado --------------------------------------
  % El arco NO va de 90 a 90-theta: theta es el angulo real en 3D, y sobre el papel la
  % direccion del estado esta en atan2(\Py,\Px). Hay que arcar hasta ESE angulo para que
  % el arco muera justo sobre la flecha.
  \pgfmathsetmacro{\Apsi}{atan2(\Py,\Px)}
  \pgfmathsetmacro{\Rt}{0.46*\R}
  \draw (0,\Rt) arc[start angle=90, end angle=\Apsi, radius=\Rt];
  \pgfmathsetmacro{\Amid}{0.5*(90+\Apsi)}
  \node at ({1.24*\Rt*cos(\Amid)},{1.24*\Rt*sin(\Amid)}) {$\theta$};

  % --- angulo phi: en el plano ecuatorial, del eje x a la proyeccion -----------------
  % Se dibuja sobre la MISMA elipse del ecuador, a radio reducido, para que se lea como
  % un angulo tumbado en el plano y no como un angulo en el plano del papel.
  \pgfmathsetmacro{\Rp}{0.42*\R}
  \draw plot[domain=\Tx:\Tp, samples=40, smooth]
        ({\Rp*cos(\x)},{\K*\Rp*sin(\x)});
  \pgfmathsetmacro{\Pmid}{\Tx+0.5*\Phi}
  \node at ({1.55*\Rp*cos(\Pmid)},{1.55*\K*\Rp*sin(\Pmid)}) {$\phi$};

\end{tikzpicture}%
\endgroup


    {\footnotesize Fuente: elaboración propia.}
\end{figure}

¿Cuánta información representa un qubit? Teóricamente, existe un número infinito de
puntos sobre la esfera unidad. Sin embargo, esta
conclusión resulta engañosa, y el motivo es el comportamiento del qubit cuando se lo observa.
Conviene recordar que medir un qubit devuelve únicamente 0 o 1. Además, la medición cambia el
estado del qubit y lo colapsa, desde su superposición de $\ket{0}$ y $\ket{1}$, al estado
concreto. Por ejemplo, si la medición de
$\ket{+}$ da 0, el estado del qubit tras la medición será $\ket{0}$.

¿Por qué se produce este colapso? Nadie lo sabe: este comportamiento es sencillamente uno de
los postulados fundamentales de la mecánica cuántica. Lo relevante para nuestros propósitos es
que de una sola medición se obtiene un único bit de información sobre el estado del qubit.
Solo midiendo infinitos qubits preparados de forma idéntica
sería posible determinar $\alpha$ y $\beta$ para un qubit en el estado dado por la
ecuación~\eqref{eq:superposicion}.

Falta una pregunta todavía: ¿cuánta información representa un qubit si
no lo medimos? La pregunta es difícil porque un estado cuántico contiene información que no podemos medir directamente. 
Sin embargo, cuando un sistema cuántico cerrado evoluciona sin realizar ninguna medición, la naturaleza 
parece ``llevar la cuenta'' de todas las variables continuas que describen el estado, como $\alpha$ y $\beta$. 
En este sentido, incluso un solo qubit encierra una gran cantidad de información oculta, y dicha cantidad 
\textbf{crece de forma exponencial con el número de qubits}.
.

\subsection{Qubits múltiples}
\label{sec:multiples}
% Apartado BREVE y con poca matematica: su unico objetivo es que se vea el crecimiento
% 2^n antes de hablar de puertas, porque despues se usa constantemente.
% Fuente: Nielsen y Chuang (2010), seccion 1.2.
% ⚠️ VERIFICAR LA PAGINA antes de la entrega. Se cita por seccion.
Un solo qubit no basta para calcular nada interesante. Entonces, ¿cómo
se describe el estado conjunto de varios? La respuesta explica por sí sola por
qué un ordenador cuántico puede hacer cosas que uno clásico no.

Con \textbf{dos qubits} las combinaciones posibles de valores clásicos son cuatro: 00, 01,
10 y 11. A cada una le corresponde un estado de la base, y el estado conjunto es una
superposición de los cuatro, con una amplitud compleja por cada uno:
\begin{equation}
\label{eq:dos_qubits}
\ket{\psi} = \alpha_{00}\ket{00} + \alpha_{01}\ket{01}
           + \alpha_{10}\ket{10} + \alpha_{11}\ket{11}
\;\equiv\;
\begin{pmatrix} \alpha_{00} \\ \alpha_{01} \\ \alpha_{10} \\ \alpha_{11}\end{pmatrix}.
\end{equation}
La notación $\ket{00}$ es una abreviatura de $\ket{0}\otimes\ket{0}$, el
\textbf{producto tensorial} de los estados de cada qubit. La operación es más simple de lo
que sugiere el símbolo: consiste en multiplicar cada componente del primer vector por todas
las del segundo y apilar los resultados, de forma que dos vectores de dos componentes
producen uno de cuatro,
\begin{equation}
\label{eq:tensorial}
\begin{pmatrix} a \\ b \end{pmatrix} \otimes \begin{pmatrix} c \\ d \end{pmatrix}
=
\begin{pmatrix} a\,c \\ a\,d \\ b\,c \\ b\,d \end{pmatrix},
\qquad\text{y en particular}\qquad
\ket{0}\otimes\ket{0} =
\begin{pmatrix} 1 \\ 0 \end{pmatrix} \otimes \begin{pmatrix} 1 \\ 0 \end{pmatrix}
=
\begin{pmatrix} 1 \\ 0 \\ 0 \\ 0 \end{pmatrix} = \ket{00}.
\end{equation}
Es también la razón de que el número de amplitudes se multiplique por dos con cada qubit
que se añade: cada componente del estado anterior se desdobla en dos. Y la condición de
normalización se mantiene: la suma de los
módulos al cuadrado de las cuatro amplitudes vale uno, porque al medir se obtiene una de
las cuatro cadenas de bits con esas probabilidades.

Con \textbf{tres qubits} el razonamiento se repite. Las cadenas posibles pasan a ser ocho
—de 000 a 111—, y el estado necesita ocho amplitudes:
\begin{equation}
\label{eq:tres_qubits}
\ket{\psi} = \alpha_{000}\ket{000} + \alpha_{001}\ket{001} + \cdots + \alpha_{111}\ket{111}
\;\equiv\;
\begin{pmatrix} \alpha_{000} & \alpha_{001} & \cdots & \alpha_{111}\end{pmatrix}^{\!\intercal}.
\end{equation}

Y así sucesivamente:
\begin{itemize}
    \item con \textbf{cuatro qubits}, los estados de la base van de $\ket{0000}$ a
          $\ket{1111}$ y hacen falta \textbf{16 amplitudes};
    \item con \textbf{cinco qubits}, van de $\ket{00000}$ a $\ket{11111}$ y hacen falta
          \textbf{32 amplitudes};
    \item \dots
\end{itemize}

El patrón está claro: \textbf{$n$ qubits requieren $2^n$ amplitudes complejas}. Con 50
qubits son más de mil billones de números, y describir ese estado en un ordenador clásico
es imposible. Ahí está la promesa de la computación cuántica.

\subsection{Puertas cuánticas y circuitos}
\label{sec:puertas}
% Fuente de los dos primeros parrafos: Nielsen y Chuang (2010), seccion 1.3.
% ⚠️ VERIFICAR LA PAGINA antes de la entrega. Se cita por seccion porque no esta
%    comprobada; no inventarla.
%
% ⚠️ Del original se ha DESCARTADO la promesa «including a circuit which teleports
%    qubits!»: anuncia un contenido que esta memoria no va a desarrollar, y prometer
%    algo que luego no aparece es de las cosas que un tribunal si nota.
Los cambios que experimenta un estado cuántico se pueden describir con el lenguaje de la
computación cuántica. Del mismo modo que un ordenador clásico se construye a partir de un
circuito eléctrico formado por cables y puertas lógicas, un ordenador cuántico se construye
a partir de un circuito cuántico formado por cables y puertas cuánticas elementales,
encargadas de transportar y manipular la información cuántica
\cite[secc.~1.3]{nielsen2010quantum}.

Los circuitos de un ordenador clásico constan de cables y puertas lógicas. Los cables
sirven para transportar la información a lo largo del circuito, mientras que las puertas
lógicas realizan manipulaciones sobre esa información y la convierten de una forma en otra.

% La forma vectorial ya se introdujo en el apartado 2.1.1; aqui solo se recupera.
Escrita en forma vectorial —ecuación~\eqref{eq:vector_estado}—, la información de un qubit
es una columna de dos números complejos. Escrito así, el papel de una puerta cuántica se deduce solo: \textbf{una puerta es una
matriz que transforma el vector de estado en otro}. Si $U$ denota la puerta y $\ket{\psi}$
el estado de partida, el estado resultante es $\ket{\psi'} = U\ket{\psi}$, y aplicar varias
puertas en secuencia equivale a multiplicar sus matrices.

\paragraph{Un ejemplo.} La puerta $X$, análoga cuántica de la negación clásica, se define
como
\begin{equation}
\label{eq:puerta_x}
X = \begin{pmatrix} 0 & 1 \\ 1 & 0 \end{pmatrix},
\qquad\text{de modo que}\qquad
X\begin{pmatrix} \alpha \\ \beta \end{pmatrix} = \begin{pmatrix} \beta \\ \alpha \end{pmatrix}.
\end{equation}
Sobre los estados de la base actúa como cabría esperar de una negación, llevando $\ket{0}$
a $\ket{1}$ y $\ket{1}$ a $\ket{0}$. Sobre una superposición, en cambio, hace algo que no
tiene equivalente clásico: intercambia las dos amplitudes sin necesidad de medir, y por
tanto sin destruir la superposición.

% ⚠️ FORMATO: titulo arriba y fuente abajo (Instrucciones §1.3, p. 6).
\begin{table}[htbp]
    \centering
    \caption{Puertas cuánticas de un qubit más habituales.}
    \label{tab:puertas}
    \begin{tabular}{@{}llp{0.46\textwidth}@{}}
        \toprule
        \textbf{Puerta} & \textbf{Matriz} & \textbf{Efecto} \\
        \midrule
        $X$ & $\begin{pmatrix} 0 & 1 \\ 1 & 0\end{pmatrix}$
            & Intercambia $\ket{0}$ y $\ket{1}$. Es la negación cuántica. \\[1.1em]
        $Y$ & $\begin{pmatrix} 0 & -i \\ i & 0\end{pmatrix}$
            & Intercambia los estados e introduce además una fase. \\[1.1em]
        $Z$ & $\begin{pmatrix} 1 & 0 \\ 0 & -1\end{pmatrix}$
            & Deja $\ket{0}$ intacto y cambia el signo de $\ket{1}$. \\[1.1em]
        $H$ & $\dfrac{1}{\sqrt{2}}\begin{pmatrix} 1 & 1 \\ 1 & -1\end{pmatrix}$
            & Puerta de Hadamard: crea superposición a partir de un estado de la base. \\[1.1em]
        $S$ & $\begin{pmatrix} 1 & 0 \\ 0 & i\end{pmatrix}$
            & Aplica una fase de $\pi/2$ al estado $\ket{1}$. \\[1.1em]
        $T$ & $\begin{pmatrix} 1 & 0 \\ 0 & e^{i\pi/4}\end{pmatrix}$
            & Aplica una fase de $\pi/4$. \\[1.1em]
        $R_z(\theta)$ & $\begin{pmatrix} 1 & 0 \\ 0 & e^{i\theta}\end{pmatrix}$
            & Rotación de fase de ángulo arbitrario en torno al eje $z$. \\
        \bottomrule
    \end{tabular}

    {\footnotesize Fuente: elaboración propia.}
\end{table}

Como toda puerta de un qubit es unitaria, su efecto sobre la esfera de Bloch de la
Figura~\ref{fig:bloch} es siempre una \textbf{rotación}: el estado se desplaza sobre la
superficie. La Figura~\ref{fig:hadamard} es un ejemplo de como la puerta de
Hadamard, que puede descomponerse en un giro de $90^\circ$ en torno al eje $y$ seguido de
otro de $180^\circ$ en torno al eje $x$, actúa sobre el estado
$(\ket{0}+\ket{1})/\sqrt{2}$ y lo lleva de vuelta a $\ket{0}$.

% ⚠️ FORMATO: titulo arriba y fuente abajo (Instrucciones §1.3, p. 6).
% ⚠️ FUENTE: es un redibujo de una figura de Nielsen y Chuang, no una invencion propia.
%    La forma academica correcta es «elaboracion propia a partir de…», no «elaboracion
%    propia» a secas. Verificar el numero de figura del libro y añadirlo si se quiere.
% Codigo de color: azul = el qubit · naranja = la transformacion.
\begin{figure}[htbp]
    \centering
    \caption{Acción de la puerta de Hadamard sobre la esfera de Bloch.}
    \label{fig:hadamard}
    % =====================================================================================
% La puerta de Hadamard vista sobre la esfera de Bloch — figura propia en TikZ.
%
% Tres viñetas: el estado de partida (|0>+|1>)/raiz(2) sobre el eje x, la rotacion de
% 90 grados en torno a y que lo lleva al polo sur, y la rotacion de 180 grados en torno
% a x que lo sube al polo norte. Resultado: H|+> = |0>.
%
% CODIGO DE COLOR, deliberado y constante en toda la memoria:
%   azulunir   -> el QUBIT (vector de estado, punto y etiqueta)
%   acentofig  -> las TRANSFORMACIONES (giros, planos, flechas de operacion)
%
% ⚠️ \shorthandoff{<>} es obligatorio: `babel` en español hace activos < y >, y eso
%    rompe la sintaxis de puntas de flecha de TikZ.
% =====================================================================================
\begingroup
\shorthandoff{<>}%
\begin{tikzpicture}[>={Stealth[length=2mm]}, line width=0.5pt, scale=0.92]

  \pgfmathsetmacro{\R}{1.45}    % radio de cada esfera
  \pgfmathsetmacro{\K}{0.30}    % achatamiento del ecuador; ver nota al pie del fichero
  \pgfmathsetmacro{\Tx}{235}    % angulo de pantalla del eje x
  \pgfmathsetmacro{\D}{4.9}     % separacion entre viñetas

  % --- una esfera con sus tres ejes, reutilizable ------------------------------------
  \newcommand{\esferabase}{%
    \draw (0,0) circle (\R);
    \draw[dashed] (0,0) ellipse ({\R} and {\K*\R});
    \draw[->] (0,0) -- (0,\R);
    \node[right=1pt] at ({0.02*\R},{0.78*\R}) {\scriptsize $\hat{\mathbf{z}}$};
    \draw[->] (0,0) -- (\R,0) node[right=-1pt] {\scriptsize $\hat{\mathbf{y}}$};
    \draw[->] (0,0) -- ({\R*cos(\Tx)},{\K*\R*sin(\Tx)})
          node[below left=-3pt] {\scriptsize $\hat{\mathbf{x}}$};%
  }

  % =====================================================================================
  % 1) Estado de partida: (|0>+|1>)/raiz(2), sobre el eje x
  % =====================================================================================
  \begin{scope}[xshift=0cm]
    \esferabase
    \draw[azulunir, line width=1.1pt]
          (0,0) -- ({\R*cos(\Tx)},{\K*\R*sin(\Tx)});
    \fill[azulunir] ({\R*cos(\Tx)},{\K*\R*sin(\Tx)}) circle (2pt);
    \node[azulunir, anchor=east] at ({-1.06*\R},{-0.40*\R})
         {\small $(\ket{0}+\ket{1})/\sqrt{2}$};
    % giro de 90 grados en torno al eje y
    \draw[acentofig, ->, line width=0.7pt]
          ({0.30*\R},{-0.038*\R}) arc[start angle=-170, end angle=110,
          x radius={0.14*\R}, y radius={0.22*\R}];
  \end{scope}

  \draw[->, line width=0.8pt] ({0.5*\D-0.45},0) -- ({0.5*\D+0.45},0);

  % =====================================================================================
  % 2) Tras el giro: |1> en el polo sur. La rotacion de 180 grados en torno a x se
  %    representa con el plano y la flecha doble.
  % =====================================================================================
  \begin{scope}[xshift=\D cm]
    % el plano va PRIMERO: asi la esfera y el estado quedan por encima y se ven
    \draw[acentofig, fill=acentofig, fill opacity=0.06, line width=0.6pt]
          ({-1.78*\R},{-0.33*\R}) -- ({-0.90*\R},{0.31*\R})
       -- ({ 1.78*\R},{ 0.31*\R}) -- ({ 0.90*\R},{-0.33*\R}) -- cycle;
    \esferabase
    \draw[azulunir, line width=1.1pt] (0,0) -- (0,-\R);
    \fill[azulunir] (0,-\R) circle (2pt);
    \node[azulunir, anchor=north] at (0,{-1.08*\R}) {\small $\ket{1}$};
    % flecha doble: el estado atraviesa el plano de arriba abajo
    \draw[acentofig, <->, line width=0.9pt]
          ({0.58*\R},{-0.52*\R}) -- ({0.58*\R},{0.52*\R});
  \end{scope}

  \draw[->, line width=0.8pt] ({1.5*\D-0.45},0) -- ({1.5*\D+0.45},0);

  % =====================================================================================
  % 3) Resultado: |0> en el polo norte
  % =====================================================================================
  \pgfmathsetmacro{\DD}{2*\D}
    \begin{scope}[xshift=\DD cm]
    \esferabase
    \draw[azulunir, line width=1.1pt] (0,0) -- (0,\R);
    \fill[azulunir] (0,\R) circle (2pt);
    \node[azulunir, anchor=south east] at ({-0.04*\R},{1.06*\R}) {\small $\ket{0}$};
  \end{scope}

\end{tikzpicture}%
\endgroup

% =====================================================================================
% NOTA SOBRE \K, EL ACHATAMIENTO
%
% Determina desde que altura se mira la esfera, y con ello el angulo VISUAL entre los
% ejes x e y. Con el eje y horizontal, ese angulo no puede llegar a 90 grados: el eje x
% ha de apoyarse sobre la elipse del ecuador —porque el ecuador ES el plano xy— y la
% unica forma de que forme 90 grados con y seria que apuntase recto hacia abajo, encima
% del eje z. Lo medido:
%
%     K = 0,30  ->  157 grados   punta de x a 0,62·R del centro
%     K = 0,55  ->  142 grados   punta de x a 0,73·R del centro   <- el elegido
%     K = 0,75  ->  133 grados   punta de x a 0,84·R del centro
%     K = 1,00  ->  125 grados   punta de x a 1,00·R  (el ecuador se confunde con
%                                el contorno y la esfera deja de parecer una esfera)
%
% Se probo 0,55 y se descarto: el usuario prefiere la perspectiva original de 0,30.
% La tabla se deja aqui por si en algun momento se quiere volver a mover.
% =====================================================================================


    {\footnotesize Fuente: elaboración propia a partir de \citeNP{nielsen2010quantum}.}
\end{figure}

\paragraph{Ejemplo de puerta de dos qubits.} El estado conjunto de dos qubits necesita cuatro
amplitudes, una por cada combinación de la base $\ket{00}$, $\ket{01}$, $\ket{10}$ y
$\ket{11}$, de modo que una puerta que actúe sobre ambos es una matriz unitaria de
$4\times 4$. La más conocida es la \textbf{CNOT}, que niega el segundo qubit —el
\textit{objetivo}— solo cuando el primero —el \textit{control}— vale 1:
\begin{equation}
\label{eq:cnot}
\mathrm{CNOT} =
\begin{pmatrix}
1 & 0 & 0 & 0 \\
0 & 1 & 0 & 0 \\
0 & 0 & 0 & 1 \\
0 & 0 & 1 & 0
\end{pmatrix}.
\end{equation}

\paragraph{El circuito.} Un circuito cuántico es una secuencia de puertas aplicadas sobre
un conjunto de qubits, seguida de una medición. Se representa con una línea horizontal por
qubit —los «cables» del símil clásico— y las puertas colocadas sobre ellas en el orden en
que se ejecutan. Al circuito lo describen dos magnitudes: la \textbf{anchura},
que es el número de qubits, y la \textbf{profundidad}, que es el número de capas de puertas
que deben ejecutarse una tras otra.

La Figura~\ref{fig:circuito_ghz} muestra un ejemplo concreto. Los cinco cables horizontales
son los qubits; sobre el primero actúa una puerta de Hadamard y a continuación una cadena de
cuatro CNOT propaga el entrelazamiento al resto. El circuito tiene por tanto anchura cinco y
profundidad cinco, y prepara el estado
$(\ket{00000}+\ket{11111})/\sqrt{2}$, en el que los cinco qubits están entrelazados entre sí.

% ⚠️ FORMATO: titulo arriba y fuente abajo (Instrucciones §1.3, p. 6).
% Figura generada por el propio proyecto: TFM-Quantum/figures/00_anatomia/circuitos/.
% Se copia el PDF, no el PNG, porque es vectorial y no pixela.
%
% ⚠️ La imagen trae su PROPIO titulo incrustado («GHZ — Estado GHZ — el entrelazamiento
%    en su forma mas pura»), asi que se duplica con el \caption. Lo suyo es regenerarla
%    sin titulo desde src/viz/circuitos, o recortarlo. Se deja asi de momento.
\begin{figure}[htbp]
    \centering
    \caption{Tres ejemplos de circuito: preparación de un estado GHZ, un ansatz
             \textit{hardware-efficient} y una transformada cuántica de Fourier.}
    \label{fig:circuito_ghz}
    % Composicion compacta: GHZ y HEA en una fila, y la QFT —mas larga— debajo,
    % ocupando el ancho completo. Sin paquetes nuevos: dos minipages y un salto.
    % ⚠️ Las tres imagenes traen SU PROPIO titulo incrustado, que chocaba con el
    %    \caption y quedaba descolgado a la izquierda. Se recortan los 20 bp
    %    superiores de cada PDF —medido: el rotulo ocupa el 6,3 % del alto—.
    %    Lo limpio seria regenerarlas sin titulo desde TFM-Quantum/src/viz/circuitos.
    \includegraphics[trim=0 0 0 20bp, clip, width=0.47\textwidth]{cap02/circuito_ghz}\hfill
    \includegraphics[trim=0 0 0 20bp, clip, width=0.47\textwidth]{cap02/circuito_hea}

    \vspace{0.4em}
    \includegraphics[trim=0 0 0 20bp, clip, width=0.88\textwidth]{cap02/circuito_qft}

    {\footnotesize Fuente: elaboración propia.}
\end{figure}

\paragraph{La base nativa.} Un procesador real no implementa todas las puertas de la
Tabla~\ref{tab:puertas}, sino un conjunto reducido a partir del cual puede construirse
cualquier otra. Antes de ejecutarse, todo circuito se \textbf{traduce} a ese conjunto
mediante un proceso llamado \textit{transpilación}. En los procesadores IBM Heron la base
nativa es $\{cz, id, rz, sx, x\}$.

Esa traducción evidencia más problema que aborda este trabajo. Un
circuito teóricamente tiene diez puertas puede convertirse en cientos al expresarse en
la base nativa, y \textbf{cada puerta ejecutada es una oportunidad más de introducir
error}. Cuanto más largo es el circuito traducido, menos señal sobrevive hasta la medición.

% =====================================================================================
% PENDIENTE en esta subseccion:
%   - FIGURA de un circuito sencillo (dos qubits, H + CNOT + medicion). Seria muy util
%     para un tribunal que no sabe de cuantica: ahora mismo el circuito se describe con
%     palabras pero no se ve. Se puede dibujar en TikZ, como la esfera de Bloch, o con
%     el paquete `quantikz` si se decide añadirlo al preambulo.
%   - Verificar la pagina exacta de la cita de Nielsen y Chuang.
% =====================================================================================

\subsection{El circuito como grafo acíclico dirigido}
\label{sec:dag}
% ⚠️ ESTA SUBSECCION ERA LA BISAGRA hacia el 2.2: cerraba el bloque cuantico con la idea
%    de que, si un circuito es un grafo, predecirlo es regresion sobre grafos. Al subirla
%    aqui quedan DOS apartados por debajo (hardware y utilidad), asi que la transicion ya
%    no cae junto al 2.2. Hay que reponerla al final de «Para que sirve un ordenador
%    cuantico», o el 2.2 arrancara en seco.
Un circuito \textbf{es} un DAG: cada puerta es un nodo, las dependencias son las aristas.
Profundidad como camino más largo.

% ⚠️ FORMATO: titulo arriba y fuente abajo (Instrucciones §1.3, p. 6).
% Ambas figuras las genera figuras/cap02/gen_circuito_dag.py, que las produce con el
% propio Qiskit: el grafo NO es una ilustracion inventada, es la representacion interna
% que el transpilador manipula de verdad. El coloreado de los nodos —verde entradas,
% azul operaciones, rojo salidas— lo pone `dag_drawer` por si mismo.
\begin{figure}[htbp]
    \centering
    \caption{Un circuito y su grafo acíclico dirigido: cada puerta es un nodo y cada
             arista, la dependencia sobre un qubit.}
    \label{fig:dag}
    \begin{minipage}[c]{0.44\textwidth}
        \centering
        \includegraphics[width=\textwidth]{cap02/dag_circuito}
    \end{minipage}\hfill
    \begin{minipage}[c]{0.52\textwidth}
        \centering
        \includegraphics[width=0.92\textwidth]{cap02/dag_grafo}
    \end{minipage}

    {\footnotesize Fuente: elaboración propia.}
\end{figure}

\textbf{Y aquí es donde la física se convierte en un problema de aprendizaje automático:} si
un circuito es un grafo, predecir su comportamiento es un problema de regresión sobre
grafos. Eso es lo que motiva la sección siguiente.


% =====================================================================================
% =====================================================================================
% 2.1.6 — EL HARDWARE ACTUAL
% Sustituye al bloque de comentarios que habia bajo \subsection{El hardware actual}.
% Enfoque acordado: divulgativo («a modo curiosidad»), sin tecnicismos, y quedandose
% solo con los datos que el resto de la memoria necesita: la topologia limitada, la
% recalibracion diaria y el coste de simular. Las cifras van fechadas en 2026 y con cita.
% =====================================================================================

\subsection{El hardware actual}
\label{sec:hardware}

Conviene hacerse una idea de qué son físicamente las máquinas de las que habla esta
memoria, aunque solo sea para situar con una idea general el problema. Un procesador cuántico
superconductor como los de IBM es un chip del tamaño de una moneda que vive en el fondo de
un criostato: un cilindro de cobre dorado, de casi dos metros, cuya única función es
mantener el chip a unas centésimas de grado por encima del cero absoluto. Los qubits no se
programan escribiendo en una memoria, sino disparándoles pulsos de microondas por unos
cables coaxiales; cada puerta del circuito es un pulso. Y la medición no devuelve un
resultado, sino una estadística.

% Parrafo nuevo: explica QUE significa «repetir miles de veces y contar». Es el que
% justifica que el objetivo del TFM sea un valor esperado y no una salida determinista.
Nos detenemos en esto último. Al medir un qubit no
se lee el estado en el que estaba: el estado \textbf{se destruye} y en su lugar aparece un
cero o un uno, elegido al azar con la probabilidad que marcaba ese estado. Una única
ejecución, por tanto, no informa de nada; devuelve una cadena de ceros y unos que podría
haber salido de otra máquina completamente distinta. Lo que se hace en la práctica es
ejecutar el mismo circuito miles de veces seguidas —cada repetición se llama
\emph{shot}— y anotar cuántas veces sale cada resultado. De ese recuento se obtiene el
promedio que interesaba, y el circuito se ha ejecutado entero, desde cero, en cada una de
esas miles de repeticiones. La consecuencia se todo esto es que la precisión es un factor que se paga caro, ya que para reducir el error
en los resultados es necesario ejecutar más veces los circuitos, lo que aumenta coste de uso del procesador.

Merece de interés mencionar dos características del hardware que lo hacen distinto de un ordenador clásico:

\textbf{Primera: los qubits no se hablan todos con todos.} En un chip clásico cualquier
dato puede llegar a cualquier registro; en uno cuántico cada qubit está conectado solo a
dos o tres vecinos. Si un algoritmo necesita que interactúen dos qubits que no lo son, la
única salida es ir acercando el estado de uno hasta el otro con puertas intermedias, y cada
puerta añadida es ruido añadido.

\textbf{Segunda: el chip no es el mismo dos días seguidos.} Estas máquinas se recalibran a
diario, y cada calibración publica un boletín con la tasa de error de cada puerta y los
tiempos de vida de cada qubit. Esos números cambian y un qubit que ayer era de los
buenos hoy puede ser de los malos, y hay puertas que se estropean y no se reparan.
Esta es la razón de que el conjunto de datos del Capítulo~\ref{cap:dataset}
contenga la calibración real de un intervalo de días distintos en lugar de una única
fotografía del chip.

% ⚠️ FORMATO: titulo arriba y fuente abajo (Instrucciones §1.3, p. 6).
% Tabla simplificada a proposito: solo el catalogo de chips. El comentario sobre el
% cambio de estrategia (subir a 1121 y bajar a 156) se retiro por peticion expresa.
\begin{table}[htbp]
    \centering
    \caption{Familias de procesadores cuánticos de IBM y su número de qubits.}
    \label{tab:familias}
    \begin{tabular}{@{}lrr@{}}
        \toprule
        \textbf{Familia} & \textbf{Año} & \textbf{Qubits} \\
        \midrule
        Falcon    & 2019 & 27 \\
        Eagle     & 2021 & 127 \\
        Osprey    & 2022 & 433 \\
        Condor    & 2023 & 1\,121 \\
        Heron     & 2023 & 133 / 156 \\
        Nighthawk & 2025 & 120 \\
        \bottomrule
    \end{tabular}

    {\footnotesize Fuente: elaboración propia a partir de \citeNP{ibm2026hardware}.}
\end{table}

% Bloque nuevo: la equivalencia chip cuantico <-> ordenador clasico.
% Cifras calculadas como 2^n amplitudes x 16 bytes (complejo de doble precision).
% El 17,2 GB de n=15 es dato del propio proyecto: con ruido el coste es 4^n = 2^30.

Aunque 156 qubits parecen poca cosa habría que repasar \textbf{cuánta memoria haría falta en un ordenador normal
para imitar el chip}. Y la respuesta es que describir el estado de $n$ qubits exige guardar
$2^n$ números, de modo que \textbf{cada qubit que se añade duplica} el coste. 
\begin{itemize}
    \item Con 30 qubits hacen falta unos 17 GB, la memoria de un portátil actual bueno.
    \item  Con 40 qubits, unos 17 TB, que ya es un servidor grande.
    \item Con 50 qubits, unos 18 PB: un centro de datos entero dedicado a imitar un chip del tamaño de una moneda. 
    \item Y con los 156 qubits de un Heron sale un número de 48 cifras que ninguna máquina actual ni futura podrá almacenar.
\end{itemize}

Estas ejecuciones con ruido de chips cuánticos se pueden simular con ordenadores clásicos mediante paquetes de 
simuladores, como Qiskit Aer. Debido al gran cómputo que exigen una ejecución cuántica, los 
circuitos simulados en este trabajo llegan hasta 15 qubits.

Con todo esto hemos visto que el chip real maneja 156 qubits sin problema y el ordenador clásico no es capaz
ni de acercarse. Por tanto, no se puede calcular de antemano cuánto va a degradar
el ruido un resultado, porque ese cálculo es precisamente lo que resulta imposible. Lo
que sí puede intentarse es \textbf{aprenderlo} a partir de ejemplos.

\section{Fundamentos de los modelos de aprendizaje automático}
\label{sec:fund_ia}
% =====================================================================================
% SOLO TEORIA. Que es cada familia de modelos y como funciona. Ni una cifra nuestra,
% ni una decision de diseño: eso es del capitulo 9.
%
% ⚠️ REFERENCIAS PENDIENTES DE VERIFICAR antes de citarlas aqui (no estan en el .bib):
%      - Ridge original: Hoerl y Kennard (1970), Technometrics
%      - Random Forest: Breiman (2001), Machine Learning
%      - Atencion / Transformer: Vaswani y cols. (2017), NeurIPS
%      - Un manual general de referencia (Hastie/Tibshirani/Friedman o Bishop)
%    ➡️ Verificar DOI de cada una. Regla del proyecto: ninguna cita sin DOI comprobado.

\subsection{Dos formas de representar un circuito}
\label{sec:dos_representaciones}
% APERTURA DE 2.2. Cinco lineas, pero son las que convierten las tres subsecciones
% siguientes en una comparativa en marcha y no en tres fichas sueltas.
Del grafo del apartado anterior se puede partir de dos maneras. La primera es
\textbf{resumirlo en una fila de números} —cuántas puertas, de qué tipo, cuánto dura— y
aplicar un modelo tabular clásico. La segunda es \textbf{dárselo al modelo tal cual}, con su
estructura intacta.

Esa disyuntiva es la que estructura todo el trabajo, y las tres familias que siguen son las
candidatas: dos que exigen la tabla y una que acepta el grafo.

\subsection{Regresión lineal regularizada: Ridge}
\label{sec:teoria_ridge}
Mínimos cuadrados ordinarios y qué optimizan. El problema de la \textbf{colinealidad}: cuando
dos columnas dicen casi lo mismo, la solución se vuelve inestable y los coeficientes se
disparan.

La \textbf{penalización $L_2$} y el compromiso sesgo--varianza: se acepta algo de sesgo a
cambio de coeficientes acotados. Interpretación geométrica.

\subsection{Ensembles de árboles: Random Forest}
\label{sec:teoria_rf}
El árbol de decisión como partición recursiva del espacio de entrada. Su problema: varianza
alta.

\textbf{Bagging} y aleatorización de variables en cada división: por qué promediar muchos
árboles decorrelacionados reduce la varianza sin aumentar apenas el sesgo.

\textbf{Una propiedad teórica que importará después:} la predicción es siempre un promedio de
valores vistos en el entrenamiento, así que el modelo \textbf{no puede extrapolar} fuera del
rango aprendido. Aquí solo se enuncia; sus consecuencias se miden en el
capítulo~\ref{cap:comparativa}.

\subsection{Redes neuronales sobre grafos}
\label{sec:teoria_gnn}
% AQUI SI hay que detenerse: es lo que el tribunal probablemente no ha visto aplicado
% a este problema.
%
% 🔴 CORREGIDO (sep-2026). Este apartado decia lo CONTRARIO de la decision vigente:
%    planteaba la atencion como alternativa al paso de mensajes. La arquitectura del
%    trabajo es un MPNN CON paso de mensajes y nodo virtual BIDIRECCIONAL; el Graph
%    Transformer se descarto, y con medidas. Ver doc/00_Justificaciones.md §5.
%
% ⚠️ AQUI SOLO VA LA TEORIA (que es una GNN, que es el paso de mensajes, que es un nodo
%    virtual). El POR QUE de la eleccion —la matriz de estructura constante en el 99 %,
%    las distancias de 30-70 pasos— es una decision propia y medida: va en el
%    capitulo~\ref{cap:comparativa}, no aqui.
Qué es una red neuronal sobre grafos y qué la distingue de una red sobre vectores o
imágenes.

El \textbf{paso de mensajes}: cada nodo agrega información de sus vecinos, y hacen falta
tantas capas como distancia se quiera cubrir.

\textbf{El nodo virtual} como forma de obtener una representación de todo el grafo: un nodo
ficticio conectado a los demás, cuyo estado final resume la estructura entera.

% PENDIENTE de redactar: mencionar tambien la ALTERNATIVA basada en atencion, aunque sea
% brevemente, porque el capitulo 4 cita GTraQEM y el 9 explica por que se descarto. Sin
% saber que existe, ese descarte no se entiende.
